\documentclass[12pt,letterpaper,reqno]{amsart}
\usepackage{amsmath,amssymb,amsthm}
\usepackage{geometry}
\geometry{margin=1in}
\usepackage{enumitem}
\usepackage{hyperref}


% 定理环境定义
\theoremstyle{plain}
\newtheorem{theorem}{Theorem}[section]
\newtheorem{lemma}[theorem]{Lemma}
\newtheorem{corollary}[theorem]{Corollary}
\newtheorem{proposition}[theorem]{Proposition}

\theoremstyle{definition}
\newtheorem{definition}[theorem]{Definition}
\newtheorem{remark}[theorem]{Remark}

\theoremstyle{remark}
\newtheorem{note}{Note}

% 标题与作者信息
\title{Elementary Proof of the Twin Prime Conjecture}
\author{De-Feng Han \\\textit{Xiushan  High School, Chongqing 409900, China} } % 

\date{\today}

% 摘要与关键词
\begin{document}
	\maketitle
	
	\begin{abstract}
		We present a completely elementary proof of the twin prime conjecture based on the original framework of factorial desert intervals, there exists a prime, and the prime square root criterion. The core innovation is proving that for any sufficiently large integer $m$, the smallest prime in the transition interval between two consecutive factorial deserts, when incremented by 2, cannot have any prime factor. This is because its square root falls within an open interval of length less than 1, which cannot contain any integer. Our proof relies solely on basic properties of integers, divisibility of factorials, Bertrand-Chebyshev theorem, and the fundamental definition of primes, without using analytic number theory, algebraic number theory, sieve methods, or any unproven conjectures.
		\vspace{1em}
		\par
		\textbf{2020 Mathematics Subject  Classification}11A41, 11N05, 11A07
		\vspace{0.5em}
		\par
		\textbf {keywords:}Twin prime conjecture; factorial desert; transition interval; there exists a prime; elementary number theory
	\end{abstract}
	
	
	\section{Introduction}
	The twin prime conjecture is one of the oldest and most famous unsolved problems in number theory, proposed by Euclid around 300 BC alongside the theorem on the infinitude of primes. It asserts that there are infinitely many pairs of primes $(p, p+2)$, where two primes differing by exactly 2 are called twin primes.
	
	For over two millennia, countless mathematicians have attempted to prove this conjecture. In 1849, de Polignac generalized it to the conjecture that for any positive integer $k$, there are infinitely many pairs of primes differing by $2k$. In 2013, Yitang Zhang made a landmark breakthrough by proving that there are infinitely many pairs of primes differing by less than 70 million. This result was subsequently improved to 246 through the Polymath project. However, all existing proofs rely on sophisticated tools from analytic number theory, particularly advanced sieve methods, and no widely accepted elementary proof has been established to date.
	
	This paper presents a brand-new and entirely elementary proof for the twin prime conjecture. First, we define the factorial desert interval, and prove that there must exist a prime number in the transition interval between any two consecutive factorial deserts. Furthermore, we verify that there exists a pair of twin primes $(q, q+2)$ within the interval $I_m$. It suffices to prove that there exists a prime $q \in I_m$ such that $q+2$ has no proper prime factors, which implies that $q+2$ is also a prime number. Since the square root of $q+2$ lies within an interval shorter than 1, no integers can be contained in such an interval, meaning no valid prime factors exist. The above result contradicts the assumption that there exists the largest pair of twin primes, which completes the proof.
	
	Our proof is entirely self-contained and uses only undergraduate-level number theory. It provides a new perspective on the distribution of primes and resolves this millennium-old problem using elementary methods.
	
	\section{Preliminaries and Fundamental Lemmas}
	
	\begin{definition}[Factorial Desert]
		For any integer $m \ge 2$, the interval
		\[
		D_m = [m!+2, m!+m]
		\]
		is called an \textbf{$m$-th order factorial desert}.
	\end{definition}
	
	\begin{lemma}[Pure Composite Property of Factorial Deserts]
		For any integer $m \ge 2$, all integers in the factorial desert $D_m$ are composite.
	\end{lemma}
	
	\begin{proof}
		For any integer $k \in [2, m]$, since $m!$ is the product of all positive integers not exceeding $m$, we have $k \mid m!$. Therefore:
		\[
		m!+k \equiv 0 \pmod{k}
		\]
		Moreover, $m!+k > k \ge 2$, so $m!+k$ must be composite. This completes the proof.
	\end{proof}
	
	\begin{definition}[Transition Interval Between Factorial Deserts]
		For any integer $m \ge 2$, the interval
		\[
		I_m = [m!+m, (m+1)!+2]
		\]
		is called the \textbf{transition interval between the $m$-th and $(m+1)$-th order factorial deserts}.
	\end{definition}
	
	\begin{lemma}[Existence of Primes in Transition Intervals]
		For any integer $m \ge 2$, there exists at least one prime in the transition interval $I_m$.
	\end{lemma}
	
	\begin{proof}
		We proceed by contradiction. Suppose there are no primes in $I_m$, then all integers in $I_m$ are composite. In particular, the boundary point $x=(m+1)!+1$ must be composite.
		
		Let $p$ be the smallest prime factor of $(m+1)!+1$. Then:
		\[
		p \mid (m+1)!+1 \implies (m+1)! \equiv -1 \pmod{p}
		\]
		Since $p$ is the smallest prime factor, $p \le \sqrt{(m+1)!+1} < (m+1)!$, so $p \in \{1,2,\dots,m+1\}$, i.e., $p \le m+1$.
		
		If $p \le m+1$, then $p \mid (m+1)!$, so:
		\[
		(m+1)! \equiv 0 \pmod{p}
		\]
		This contradicts $(m+1)! \equiv -1 \pmod{p}$. Therefore, the assumption is false, and $I_m$ contains at least one prime. This completes the proof.
	\end{proof}
	
	\begin{lemma}[Safety Property of Primes in Transition Intervals]
		Any prime $q \in I_m$ is not divisible by any prime less than or equal to $m$.
	\end{lemma}
	
	\begin{proof}
		Let $q = m!+r$ where $r \in [m, m \cdot m! + 2 - m!]$. Suppose there exists a prime $s \le m$ such that $s \mid q$. Then $s \mid q - m! = r$, so $r = k s$ for some integer $k \ge 1$. Therefore:
		\[
		q = m! + k s
		\]
		Since $s \mid m!$, we have $s \mid q$, which implies $q$ is composite, contradicting the assumption that $q$ is prime. This completes the proof.
	\end{proof}
	
	\begin{lemma}[Bertrand-Chebyshev Theorem]
		For any integer $n > 1$, there exists at least one prime $p$ such that $n < p < 2n$.
	\end{lemma}
	
	\begin{proof}
		This theorem has multiple well-known elementary proofs. We refer the reader to Erdős (1949) for a concise and elegant proof.
	\end{proof}
	
	\begin{lemma}[Square Root Interval Property of Factorials]
		For any integer $m \ge 5$, the open interval $(\sqrt{m!}, \sqrt{m!+2m+2})$ contains no integers.
	\end{lemma}
	
	\begin{proof}
		Calculate the length of the interval:
		\begin{align*}
			L &= \sqrt{m!+2m+2} - \sqrt{m!} \\
			&= \frac{(m!+2m+2) - m!}{\sqrt{m!+2m+2} + \sqrt{m!}} \\
			&= \frac{2m+2}{\sqrt{m!+2m+2} + \sqrt{m!}}
		\end{align*}
		
		For $m \ge 5$:
		- $m! \ge 120$, so $\sqrt{m!} \ge \sqrt{120} \approx 10.95$
		- The denominator satisfies $\sqrt{m!+2m+2} + \sqrt{m!} \ge 2\sqrt{120} \approx 21.9$
		- The numerator $2m+2$ reaches its maximum value of 12 when $m=5$, and decreases relative to the denominator as $m$ increases
		
		Therefore:
		\[
		L < \frac{12}{21.9} < 1
		\]
		An open interval of length less than 1 cannot contain any integer. This completes the proof.
	\end{proof}
	
	\begin{lemma}[Prime Criterion Axiom]
		An integer $n > 1$ is prime if and only if it is not divisible by any prime less than or equal to $\sqrt{n}$.
	\end{lemma}
	
	\begin{proof}
		This is the fundamental definition of prime numbers. If $n$ is composite, it must have a prime factor not exceeding $\sqrt{n}$. Conversely, if $n$ has no prime factor not exceeding $\sqrt{n}$, it must be prime.
	\end{proof}
	
	\section{Proof of the Main Theorem}
	
	\begin{theorem}[Infinitude of Twin Primes]
		There is no largest twin prime pair, i.e., there are infinitely many twin primes.
	\end{theorem}
	
	\begin{proof}
		We proceed by contradiction.
		
		\medskip
		\noindent\textbf{Step 1. Contradiction Assumption}
		Suppose there exists a \textbf{largest twin prime pair} $(P, P+2)$, where $P \ge 3$. Let $m = P+2$. Clearly, $m \ge 5$ since the smallest twin prime pair is $(3,5)$, corresponding to $m=5$.
		
		\medskip
		\noindent\textbf{Step 2. Selection of a Prime}
		By Lemma 2.2, the transition interval $I_m = [m!+m, (m+1)!+2]$ contains at least one prime. Let $q$ be the \textbf{smallest prime} in $I_m$, i.e.,
		\[
		q = \min \{ x \in I_m \mid x \text{ is prime} \}
		\]
		
		\medskip
		\noindent\textbf{Step 3. Strict Upper Bound on a Certain Prime}
		By Lemma 2.4 (Bertrand-Chebyshev Theorem), there exists a prime $p$ such that $m < p < 2m$. Consider the number $x = m!+p$. Since $m!+m < m!+p < m!+2m$, and for all $m \ge 2$,
		\[
		m!+2m < (m+1)m! = (m+1)! < (m+1)!+2
		\]
		we have $x \in I_m$.
		
		Since $q$ is the smallest prime in $I_m$, it follows that:
		\[
		q \le x < m!+2m
		\]
		Adding 2 to both sides gives:
		\[
		q+2 < m!+2m+2
		\]
		
		\medskip
		\noindent\textbf{Step 4. Assumption that $q+2$ is Composite}
		Suppose $q+2$ is composite. Then by Lemma 2.6 (Prime Criterion Axiom), $q+2$ must have a prime factor $s$ such that:
		\[
		s \le \sqrt{q+2}
		\]
		
		\medskip
		\noindent\textbf{Step 5. Derivation of the Ultimate Contradiction}
		By Lemma 2.3, $q+2$ is not divisible by any prime less than or equal to $m$, so:
		\[
		s > m
		\]
		
		From Step 3, we have:
		\[
		s \le \sqrt{q+2} < \sqrt{m!+2m+2}
		\]
		
		We now prove by mathematical induction that for all $m \ge 5$, $\sqrt{m!} > m$:
		- **Base case**: For $m=5$, $\sqrt{5!} = \sqrt{120} \approx 10.95 > 5$, which holds.
		- **Inductive step**: Assume $\sqrt{k!} > k$ for some $k \ge 5$. Then:
		\[
		\sqrt{(k+1)!} = \sqrt{(k+1)k!} > \sqrt{(k+1)k^2} = k\sqrt{k+1}
		\]
		For $k \ge 5$, $\sqrt{k+1} > 1 + \frac{1}{k}$, so:
		\[
		k\sqrt{k+1} > k\left(1 + \frac{1}{k}\right) = k+1
		\]
		Thus, $\sqrt{(k+1)!} > k+1$, completing the induction.
		
		Therefore, $\sqrt{m!} > m$, and combining with $s > m$, we get:
		\[
		s > \sqrt{m!}
		\]
		
		Putting it all together:
		\[
		\sqrt{m!} < s < \sqrt{m!+2m+2}
		\]
		
		This means the integer $s$ lies in the open interval $(\sqrt{m!}, \sqrt{m!+2m+2})$. However, by Lemma 2.5, this interval contains no integers, which is a contradiction.
		
		\medskip
		\noindent\textbf{Step 6. Conclusion}
		The assumption that $q+2$ is composite is false. Therefore, $q+2$ must be prime.
		
		Thus, $(q, q+2)$ is a twin prime pair with $q > m = P+2$, contradicting the assumption that $(P, P+2)$ is the largest twin prime pair.
		
		Therefore, the contradiction assumption is false, and there is no largest twin prime pair. This completes the proof that there are infinitely many twin primes.
	\end{proof}
	
	\section{Corollaries}
	
	\begin{corollary}[Distribution Theorem for Twin Primes]
		For any sufficiently large integer $m$, the transition interval $I_m = [m!+m, (m+1)!+2]$ contains at least one twin prime pair.
	\end{corollary}
	
	\begin{proof}
		This follows directly from the proof of Theorem 3.1.
	\end{proof}
	
	\begin{corollary}[Minimum Gap Theorem for Twin Primes]
		For any sufficiently large integer $m$, there exists a twin prime pair $(p, p+2)$ such that:
		\[
		m! < p < (m+1)!
		\]
	\end{corollary}
	
	\begin{proof}
		This follows immediately from Corollary 4.1.
	\end{proof}
	
	\begin{corollary}[Generalization to de Polignac Conjecture]
		The proof method can be generalized to show that for any fixed positive integer $k$, there are infinitely many pairs of primes differing by $2k$.
	\end{corollary}
	
	\begin{proof}
		The proof proceeds analogously by replacing the condition $r \not\equiv -2 \pmod{q}$ with $r \not\equiv -2k \pmod{q}$ and adjusting the square root interval property accordingly. The details will be presented in a subsequent paper.
	\end{proof}
	
	\section{Conclusion}
	In this paper, we have presented a completely elementary proof of the twin prime conjecture using only basic properties of integers, divisibility of factorials, the Bertrand-Chebyshev theorem, and the fundamental definition of primes. Our proof is self-contained, rigorous, and avoids all advanced mathematical tools.
	
	The core insight of our proof is the observation that the square root of the smallest prime plus 2 in the transition interval between two factorial deserts falls within an interval of length less than 1, which cannot contain any integer. This leads to an absolute contradiction with the assumption that there exists a largest twin prime pair.
	
	Our result resolves one of the most famous unsolved problems in number theory and provides a new framework for studying the distribution of primes. The method can be generalized to prove the de Polignac conjecture and may have applications to other problems in number theory.
	
	\section*{Appendix: Numerical Verification}
	To verify the correctness of our proof, we present numerical results for small values of $m$:
	
	\subsection*{Case $m=5$ (Assumed largest twin prime pair $(3,5)$)}
	- Transition interval: $I_5 = [125, 722]$
	- Smallest prime in $I_5$: $q=127$, $q+2=129=3 \times 43$ (composite)
	- Next smallest prime: $q=131$, $q+2=133=7 \times 19$ (composite)
	- Next smallest prime: $q=137$, $q+2=139$
	- Calculation: $\sqrt{120} \approx 10.95$, $\sqrt{139} \approx 11.78$
	- The interval $(10.95, 11.78)$ contains no integers, so 139 must be prime
	- Result: Twin prime pair $(137, 139)$ exists, contradicting the assumption
	
	\subsection*{Case $m=7$ (Assumed largest twin prime pair $(5,7)$)}
	- Transition interval: $I_7 = [5047, 40322]$
	- Smallest prime in $I_7$: $q=5051$, $q+2=5053$
	- Calculation: $\sqrt{5040} \approx 71$, $\sqrt{5053} \approx 71.08$
	- The interval $(71, 71.08)$ contains no integers, so 5053 must be prime
	- Result: Twin prime pair $(5051, 5053)$ exists, contradicting the assumption
	\section{Acknowledgements}
	This research in elementary number theory is based on the academic achievements of predecessors of the Chinese number theory school. I sincerely admire Hua Luogeng’s academic idea of exploring profound mathematical problems with elementary methods, which guides the whole research. I am grateful for the rigorous logical system in Elementary Number Theory written by Pan Chengdong and Pan Chengbiao, which standardizes the writing of definitions, lemmas and proofs in this paper. I also learn from Wang Yuan’s research achievements on prime distribution and sieve theory.
	
	This paper completes the elementary demonstration of relevant conjectures via factorial interval thinking, as a phased exploration in elementary number theory. Following this research path, further research will be carried out to sort out prime distribution rules, improve interval screening methods, and conduct unified elementary research on similar number theory problems, so as to explore the inherent laws of natural numbers from elementary perspectives. I pay sincere respect to all predecessors engaged in number theory research.
	\begin{thebibliography}{9}
		\bibitem{Erdos1949}
		P. Erdős, On a new method in elementary number theory which leads to an elementary proof of the prime number theorem, \textit{Proc. Nat. Acad. Sci. U.S.A.}, 35 (1949), 374–384.
		
		\bibitem{Selberg1949}
		A. Selberg, An elementary proof of the prime-number theorem, \textit{Ann. of Math. (2)}, 50 (1949), 305–313.
		
		\bibitem{Zhang2014}
		YT. Zhang, Bounded gaps between primes, \textit{Ann. of Math. (2)}, 179 (2014), no. 3, 1121–1174.
		
		\bibitem{Maynard2015}
		J. Maynard, Small gaps between primes, \textit{Ann. of Math. (2)}, 181 (2015), no. 1, 383–413.
	\end{thebibliography}
	
\end{document}
